\documentclass[10pt,a4paper]{amsart}
\usepackage[margin=19mm]{geometry}
\usepackage{amsmath,amssymb,mathtools}
\usepackage[scr=rsfs,cal=euler]{mathalfa}
\usepackage{xfrac}
\usepackage[unicode,hidelinks,bookmarks=false]{hyperref}
\newcommand{\ie}{i.e.\ }
\newcommand{\cf}{cf.\ }
\newcommand{\resp}{respectively\ }
\newcommand{\Subsection}{\subsection}
\providecommand{\nobreakdash}{\nobreak\hbox{-}}
\setlength{\emergencystretch}{2em}
\multlinegap0pt
\usepackage{float}
\usepackage{amssymb}
\usepackage{mathtools}
\usepackage[shortlabels]{enumitem}
\usepackage{dsfont}
\usepackage{tikz}
\newtheorem{theorem}{Theorem}
\usepackage{cleveref}
\crefname{theorem}{Theorem}{Theorems}
\newcommand{\N}{\mathbb{N}}
\newcommand{\R}{\mathbb{R}}
\newcommand{\br}[1]{\lb #1 \rb} % round brackets
\newcommand{\brOf}[1]{\!\br{#1}} % round brackets
\newcommand{\eqcm}{\,,} % punctuation in equations (adds space)
\newcommand{\eqfs}{\,.}
\newcommand{\lb}{\left(} %'(' Left (round) Bracket
\newcommand{\rb}{\right)} %')' Right (round) Bracket
\title{Machine-Precision Prediction of Low-Dimensional Chaotic Systems from Noise-Free Data}
\author{Christof Sch\"otz, Niklas Boers}
\date{Source: arXiv:2507.09652v2 (CC BY 4.0)}
\begin{document}
\maketitle

\noindent\emph{Source and licence.} Machine-Precision Prediction of Low-Dimensional Chaotic Systems from Noise-Free Data. Version arXiv:2507.09652v2 (CC BY 4.0), distributed by arXiv under the Creative Commons Attribution 4.0 International licence, http://creativecommons.org/licenses/by/4.0/, as recorded in the arXiv licence metadata for this exact version and shown on the abstract page https://arxiv.org/abs/2507.09652v2. Full licence notice: \texttt{licenses/arxiv-2507.09652-CC-BY-4.0.txt}.\par\smallskip

\noindent\textit{Editorial scope.} Counted unit: the article's theorem theorem (source0.tex lines 1452--1477). The article's complete original proof of it is delivered with it (source0.tex lines 1479--1564, one \texttt{proof} environment of 86 lines). No mathematics is written, repaired or re-derived: the statement and the proof are the article's own lines, in the article's own notation, and no retained line is substituted or re-flowed.  Retained uncounted as the source-local context the proof needs: lines 1342--1366.  Nothing was omitted from any retained span, so the package carries no omission record.  No retained line contains a citation, so the wrapper carries no bibliography and no bibliography entry.  No retained line contains a figure environment, an image-inclusion command or drawing code, so no image is delivered and the package declares no asset dependency.  Editorial formatting only, in addition: an AMS \texttt{amsart} preamble that restates the article's own declarations and macros used by the retained lines, namely the environments \texttt{theorem} on separate counters, together with the standard packages those lines need.  The retained environments print in this document as Theorem 1 in order of appearance, whereas the article prints the same items with the numbering of its own full document.  The complete native source of the article as distributed in the arXiv e-print for this exact version is bundled unchanged as \texttt{sources/181497/source0.tex}.
\begin{theorem}\label{thm:rk4l63:one}
	Let
	\begin{equation*}
		\begin{bmatrix}
			\zeta_{1}\\
			\zeta_{2}\\
			\zeta_{3}
		\end{bmatrix}
		:=
		\mathsf{RK4}\brOf{f_{\mathsf{L63}}, \begin{bmatrix}
				x\\
				y\\
				z
			\end{bmatrix}, h, 1}
		\eqfs
	\end{equation*}
	Define $\alpha := 2h^{-1}a^{-1} - 1$.
	Then
	\begin{align*}
		\zeta_{1} &= -\frac{h^7a^3}{192}x^2z^2(y+\alpha x) + q_1(x,y,z)\eqcm\\
		\zeta_{2} &= \frac{h^{11}a^4}{1536}x^3y^2z(y+\alpha x)^2 + q_2(x,y,z)\eqcm\\
		\zeta_{3} &= \frac{h^{11}a^4}{1536}x^3yz^2(y+\alpha x)^2 + q_3(x,y,z)\eqcm\\
	\end{align*}
	where $q_1, q_2, q_3$ are polynomials with $\deg(q_1) \leq 4$ and $\deg(q_2), \deg(q_3) \leq 7$.
\end{theorem}

\begin{theorem}
	Let $k \in \N$.
	Let
	\begin{equation*}
		\begin{bmatrix}
			\zeta_1\\
			\zeta_2\\
			\zeta_3
		\end{bmatrix}
		:=
		\mathsf{RK4}\brOf{f_{\mathsf{L63}}, \begin{bmatrix}
				x\\
				y\\
				z
			\end{bmatrix}, h, k}
		\eqfs
	\end{equation*}
	Assume $a,h\neq 0$.
	Then
	\begin{equation*}
		\deg(\zeta_1) = F_{4k+1}
		\qquad\text{and}\qquad
		\deg(\zeta_2) = \deg(\zeta_3) = F_{4k+2}
		\eqfs
	\end{equation*}
\end{theorem}

\begin{proof}
	First note that, by \cref{thm:rk4l63:one}, one RK4 step has the form
	\begin{equation*}
		\mathsf{RK4}\brOf{f_{\mathsf{L63}}, \begin{bmatrix}
				x\\
				y\\
				z
			\end{bmatrix}, h, 1} = 
		\begin{bmatrix}
			\gamma_1 x^2 y z^2 + q_1(x,y,z)\\
			\gamma_2 x^3 y^4 z + q_2(x,y,z)\\
			\gamma_3 x^3 y^3 z^2 + q_3(x,y,z)
		\end{bmatrix}
	\end{equation*}
	with $\gamma_1,\gamma_2, \gamma_3 \in\R$ and polynomials $q_1$, $q_2$, $q_3$ that cannot cancel their respective first term as they contain different monomials. As we assume $h,a\neq 0$, we have $\gamma_1,\gamma_2, \gamma_3\neq0$.
	Now, we prove the statement of the theorem by induction over $k$:
	The induction base with $k=1$ follows directly from \cref{thm:rk4l63:one} with the arguments given above:
	\begin{align*}
		\deg(\gamma_1 x^2 y   z^2 + q_1(x,y,z))  &= 5 = F_{5}\eqcm\\
		\deg(\gamma_2 x^3 y^4 z   + q_2(x,y,z))  &= 8 = F_{6}\eqcm\\
		\deg(\gamma_3 x^3 y^3 z^2 + q_3(x,y,z))  &= 8 = F_{6}\eqfs\\
	\end{align*} 
	For the induction step, let
	\begin{equation*}
		\begin{bmatrix}
			\zeta_{k,1}\\
			\zeta_{k,2}\\
			\zeta_{k,3}
		\end{bmatrix}
		:=
		\mathsf{RK4}\brOf{f_{\mathsf{L63}}, \begin{bmatrix}
				x\\
				y\\
				z
			\end{bmatrix}, h, k}
	\end{equation*}
	and assume $\deg(\zeta_{k-1,1}) = F_{4(k-1)+1}$ and $\deg(\zeta_{k-1,2}) = \deg(\zeta_{k-1,3}) = F_{4(k-1)+2}$.
	Applying \cref{thm:rk4l63:one} to the $k$-th RK4 solver step yields
	\begin{equation*}
		\mathsf{RK4}\brOf{f_{\mathsf{L63}}, \begin{bmatrix}
				x\\
				y\\
				z
			\end{bmatrix}, h, k} = 
		\mathsf{RK4}\brOf{f_{\mathsf{L63}}, \begin{bmatrix}\zeta_{k-1,1}\\\zeta_{k-1,2}\\\zeta_{k-1,3}\end{bmatrix}, h, 1} = 
		\begin{bmatrix}
			\gamma_1 \zeta_{k-1,1}^2\zeta_{k-1,2}\zeta_{k-1,3}^2 + \tilde q_1(x,y,z)\\
			\gamma_2 \zeta_{k-1,1}^3\zeta_{k-1,2}^4\zeta_{k-1,3} + \tilde q_2(x,y,z)\\
			\gamma_3 \zeta_{k-1,1}^3\zeta_{k-1,2}^3\zeta_{k-1,3}^2 + \tilde q_3(x,y,z)
		\end{bmatrix}
		\eqfs
	\end{equation*}
	The polynomials $\tilde q_1, \tilde q_2, \tilde q_3$ are of lower degree than their respective $\zeta$-terms and, thus, cannot cancel the leading monomials in the $\zeta$-terms.
	Therefore, using the induction hypothesis,
	\begin{align*}
		\deg(\zeta_{k,1}) 
		&= 
		2\deg(\zeta_{k-1,1}) + \deg(\zeta_{k-1,2}) + 2 \deg(\zeta_{k-1,3})
		\\&=
		2  F_{4(k-1)+1} + 3  F_{4(k-1)+2}
		\\&=
		2  F_{4(k-1)+3} +  F_{4(k-1)+2}
		\\&=
		F_{4(k-1)+3} +  F_{4(k-1)+4}
		\\&=
		F_{4k+1}
	\end{align*}
	and
	\begin{align*}
		\deg(\zeta_{k,2}) 
		&= 
		3\deg(\zeta_{k-1,1}) + 4\deg(\zeta_{k-1,2}) +  \deg(\zeta_{k-1,3})
		\\&=
		3  F_{4(k-1)+1} + 5 F_{4(k-1)+2}
		\\&=
		3  F_{4(k-1)+3} + 2 F_{4(k-1)+2}
		\\&=
		F_{4(k-1)+3} + 2 F_{4(k-1)+4}
		\\&=
		F_{4(k-1)+4} + F_{4(k-1)+5}
		\\&=
		F_{4k+2}
		\eqfs
	\end{align*}
	The calculation for $\deg(\zeta_{k,3}) = F_{4k+2}$ is almost the same as for $\deg(\zeta_{k,2})$.
\end{proof}

\end{document}
